% Generated from validated synthesis. Do not hand-edit.
\section*{Retrospective Signals}
\addcontentsline{toc}{section}{Retrospective Signals}
\smallskip
\noindent\textbf{推論時の計算が独立した能力軸になった} reasoning modelの登場とtest-time scaling研究が、学習時だけでなく推論時の計算配分を技術選択として前景化させた。 \hfill{\footnotesize p.~\pageref{pkg:reasoning}}\par
\smallskip
\noindent\textbf{回答から検索・操作・接続へ広がった} 検索、画面操作、ツール接続、プロトコルが、それぞれ異なる実行面として立ち上がった。 \hfill{\footnotesize p.~\pageref{pkg:action-connectivity}}\par
\smallskip
\noindent\textbf{Openは配備設計の比較へ変わった} weightsの公開だけでなくlicense、training artifacts、architecture、local/hosted deploymentを分けて読む必要が強まった。 \hfill{\footnotesize p.~\pageref{pkg:open-deployment} / p.~\pageref{pkg:efficiency-serving}}\par
\smallskip
\noindent\textbf{MultimodalityがRealtimeと生成mediaへ分岐した} 画像理解、リアルタイム音声、画像・動画生成が、別の技術層として同時に進展した。 \hfill{\footnotesize p.~\pageref{pkg:multimodal-media}}\par
\smallskip
\noindent\textbf{Executionの拡大がControlを要求した} grounding、moderation、prompt injection、evaluation、alignmentをmodel capabilityとは別のcontrol layerとして扱う必要が明確になった。 \hfill{\footnotesize p.~\pageref{pkg:control-evaluation} / p.~\pageref{pkg:synthesis}}\par
\medskip
\begin{claimboundary}[Retrospective scope]
本号は2024年後期を後日確認可能になった一次情報も用いて再構成するRetrospective Specialである。Coverage Periodと制作時点を同一視せず、vendor / project / author claimの境界は各記事とTechnical Notesで明示する。
\end{claimboundary}
\medskip
\tableofcontents
